\section{Keeping Warm: David and Abishag}

\begin{quotex} Symbolism can deceive only those who are incapable of seeing anything except the most literal meaning. They will never understand what a symbol is, because their conceptions are irremediably limited to earthly existence and the corporeal world, where, by the most naive of illusions, they try to confine all reality.
\flright{\textsc{Rene Guenon}, \emph{Man and his Becoming}}
\end{quotex}


With a little cold wind penetrating the subtropics tonight (45/7 degrees), and the expectation of finding the grounds littered with comatose iguanas in the morning, my thoughts turn tonight to the story of King David and Abishag. This is the Biblical story:

\begin{quotex}
Now King David was old and advanced in years; and although they covered him with clothes, he could not get warm. Therefore his servants said to him, ``Let a young maiden be sought for my lord the king, and let her wait upon the king, and cherish him; let her lie in your bosom, that my lord the king may be warm.'' So they sought for a beautiful maiden throughout all the territory of Israel, and found Abishag the Shunammite, and brought her to the king. The maiden was very beautiful; and she became the king's nurse and ministered to him; but the king knew her not.
\flright{\textsc{1 Kings 1:1-4}}
\end{quotex}


When a sacred text seems odd, erroneous, inconsistent, etc., it is often a signal that the text should be understood symbolically, not literally. In the story of David and Abishag, it seems odd that lying next to a woman was the only way to keep David warm. Moreover, why did it have to be with Israel's most beautiful woman when rather than any woman? That indicates that the woman had to have a special relationship to David, one that was suitable just for him.

Physics, in the theory of many worlds, shows that it is possible for multiple degrees of existence to share the same space, time, and energy. Hence, in the human state the relationship between a man and a woman can take place in three levels, even simultaneously:

\begin{itemize}


\item Gross state: two individual beings

\item Subtle state: two sharing same consciousness

\item Causal state: alchemical marriage or unity
\end{itemize}


Since Abishag was to lie next to David, presumably at night, we can assume that the story was about the dream state (or subtle state), so the presence of a corporeal maiden (even if literally true) was not the point of the story. As above, so below, hence the woman in subtle manifestation is in consciousness unlike the woman in corporeal manifestation. The dream state is related to \emph{Tejas}, the fire element. The dream state has the knowledge of internal mental objects. This is how Guenon describes it:

\begin{quotex}
The fire element is both light and heat. ... Everything relating to this touches very closely on the very nature of life, which is inseparable from heat. Aristotle's conceptions agree fully with those of the Easterners.
\end{quotex}


Here was can see that David was not just receiving physical warmth from Abishag, but, more importantly, life itself! Fire also brings light physically, but as knowledge in the subtle state. Guenon touches on this aspect:

\begin{quotex}
The dream state is intermediate between waking and deep sleep. He who knows this, advances in the path of Knowledge by his identification with Hiranyagarbha. Thus enlightened, he is in harmony with all things, for he sees the manifested Universe as the production of his own knowledge, which cannot be separated from himself.
\end{quotex}


As a reminder, Hiranyagarbha is still Non-Supreme knowledge. Hence, it is related to salvation and immortality rather than full Liberation. In a sense, then, David and Abishag are embarking on the Divine Journey through the celestial spheres. As the embodiment of an archetype, she is his Anima, a psychopomp guiding his soul through the afterlife.

Over the years, I have been logging my dreams of the Anima. The first ones were quite vague; I would wake up with the feeling that I had just met someone. However, I could remember nothing about her and usually assumed she was a living person. Over time, they carried more content. The most recent dreams are about a woman who reveals to me some deep secret. They repeat.

In one of them, she leads up a staircase to an attic. She shows me a secret room with a chest filled with manuscripts. I am allowed to read them after she leaves, but I don't remember any specific content.

In another, there is a secret room at the back of a kitchen. A short staircase leads up to a library filled with texts. Again, I am left alone to rummage through them, but awaken in total forgetfulness. The images of the attic and kitchen library are still vivid in my mind.

I used to be frustrated about forgetting the content of the mysterious texts. As I eventually learned, that is because I was expecting head knowledge, not knowledge of the heart. I wanted a theory of everything, but the dreams were telling me that I needed to change my life.

The dream will be fulfilled when the Anima stays and unites with me in consciousness. She will not leave me alone with the texts, but we will each know them in our own way. Then we begin the journey.


\flrightit{Posted on 2023-01-14 by Cologero}

\begin{center}* * *\end{center}

\begin{footnotesize}
\begin{sffamily}
\texttt{T. Freitas on 2023-01-15 at 04:05 said:}

Very interesting view of the story of David and Abishag.

\hfill

\texttt{Santiago on 2023-01-16 at 17:46 said:}

The staircases seem important. Jung has said in one of his books (I can't remember which) that staircases represent the sexual instinct (``going-up-going-down''), which seems unlikely. In your case--assuming Janet is off-duty--what is sexual about a library or a book-room?

In most of my dreams, she leads me down the staircase first. The best interpretation I can make of it comes from Dante, signifying there's more hell to see, first.

\hfill

\texttt{Kaukomieli on 2023-01-18 at 09:15 said:}

@Santiago: I believe it was Freud who made the sexual allusion since for him Everything is about the life instinct and the Death instinct, hence sex. Jung is More subtle and ``spiritual''. The staircase leading upwards would probably be interpreted by him that something is coming up into consciousness while the staircase leading downwards would mean descending into the subconscious layers of the psyche.

Guenon was mentioned in the article a couple of times. He seemed to be adamant in his convictions that Dreams cannot transmit anything spiritual. I heartily disagree with him on that issue, since as a dreamer who has an intense nocturnal life I have seen a lot of Dreams that I can only label with the term spiritual and even enlightening. Perhaps the words of Guenon have to be understood in the context of making a difference between the psychic and the spiritual.

I am beginning to sense that the whole Bible needs to be understood symbolically.


\hfill

\end{sffamily}
\end{footnotesize}

